% Options for packages loaded elsewhere
% Options for packages loaded elsewhere
\PassOptionsToPackage{unicode}{hyperref}
\PassOptionsToPackage{hyphens}{url}
%
\documentclass[
  11pt,
  letterpaper,
  DIV=11,
  numbers=noendperiod,
  twoside]{scrartcl}
\usepackage{xcolor}
\usepackage[left=1.1in, right=1in, top=0.8in, bottom=0.8in,
paperheight=9.5in, paperwidth=7in, includemp=TRUE, marginparwidth=0in,
marginparsep=0in]{geometry}
\usepackage{amsmath,amssymb}
\setcounter{secnumdepth}{3}
\usepackage{iftex}
\ifPDFTeX
  \usepackage[T1]{fontenc}
  \usepackage[utf8]{inputenc}
  \usepackage{textcomp} % provide euro and other symbols
\else % if luatex or xetex
  \usepackage{unicode-math} % this also loads fontspec
  \defaultfontfeatures{Scale=MatchLowercase}
  \defaultfontfeatures[\rmfamily]{Ligatures=TeX,Scale=1}
\fi
\usepackage{lmodern}
\ifPDFTeX\else
  % xetex/luatex font selection
  \setmathfont[]{Garamond-Math}
\fi
% Use upquote if available, for straight quotes in verbatim environments
\IfFileExists{upquote.sty}{\usepackage{upquote}}{}
\IfFileExists{microtype.sty}{% use microtype if available
  \usepackage[]{microtype}
  \UseMicrotypeSet[protrusion]{basicmath} % disable protrusion for tt fonts
}{}
\usepackage{setspace}
% Make \paragraph and \subparagraph free-standing
\makeatletter
\ifx\paragraph\undefined\else
  \let\oldparagraph\paragraph
  \renewcommand{\paragraph}{
    \@ifstar
      \xxxParagraphStar
      \xxxParagraphNoStar
  }
  \newcommand{\xxxParagraphStar}[1]{\oldparagraph*{#1}\mbox{}}
  \newcommand{\xxxParagraphNoStar}[1]{\oldparagraph{#1}\mbox{}}
\fi
\ifx\subparagraph\undefined\else
  \let\oldsubparagraph\subparagraph
  \renewcommand{\subparagraph}{
    \@ifstar
      \xxxSubParagraphStar
      \xxxSubParagraphNoStar
  }
  \newcommand{\xxxSubParagraphStar}[1]{\oldsubparagraph*{#1}\mbox{}}
  \newcommand{\xxxSubParagraphNoStar}[1]{\oldsubparagraph{#1}\mbox{}}
\fi
\makeatother


\usepackage{longtable,booktabs,array}
\newcounter{none} % for unnumbered tables
\usepackage{calc} % for calculating minipage widths
% Correct order of tables after \paragraph or \subparagraph
\usepackage{etoolbox}
\makeatletter
\patchcmd\longtable{\par}{\if@noskipsec\mbox{}\fi\par}{}{}
\makeatother
% Allow footnotes in longtable head/foot
\IfFileExists{footnotehyper.sty}{\usepackage{footnotehyper}}{\usepackage{footnote}}
\makesavenoteenv{longtable}
\usepackage{graphicx}
\makeatletter
\newsavebox\pandoc@box
\newcommand*\pandocbounded[1]{% scales image to fit in text height/width
  \sbox\pandoc@box{#1}%
  \Gscale@div\@tempa{\textheight}{\dimexpr\ht\pandoc@box+\dp\pandoc@box\relax}%
  \Gscale@div\@tempb{\linewidth}{\wd\pandoc@box}%
  \ifdim\@tempb\p@<\@tempa\p@\let\@tempa\@tempb\fi% select the smaller of both
  \ifdim\@tempa\p@<\p@\scalebox{\@tempa}{\usebox\pandoc@box}%
  \else\usebox{\pandoc@box}%
  \fi%
}
% Set default figure placement to htbp
\def\fps@figure{htbp}
\makeatother


% definitions for citeproc citations
\NewDocumentCommand\citeproctext{}{}
\NewDocumentCommand\citeproc{mm}{%
  \begingroup\def\citeproctext{#2}\cite{#1}\endgroup}
\makeatletter
 % allow citations to break across lines
 \let\@cite@ofmt\@firstofone
 % avoid brackets around text for \cite:
 \def\@biblabel#1{}
 \def\@cite#1#2{{#1\if@tempswa , #2\fi}}
\makeatother
\newlength{\cslhangindent}
\setlength{\cslhangindent}{1.5em}
\newlength{\csllabelwidth}
\setlength{\csllabelwidth}{3em}
\newenvironment{CSLReferences}[2] % #1 hanging-indent, #2 entry-spacing
 {\begin{list}{}{%
  \setlength{\itemindent}{0pt}
  \setlength{\leftmargin}{0pt}
  \setlength{\parsep}{0pt}
  % turn on hanging indent if param 1 is 1
  \ifodd #1
   \setlength{\leftmargin}{\cslhangindent}
   \setlength{\itemindent}{-1\cslhangindent}
  \fi
  % set entry spacing
  \setlength{\itemsep}{#2\baselineskip}}}
 {\end{list}}
\usepackage{calc}
\newcommand{\CSLBlock}[1]{\hfill\break\parbox[t]{\linewidth}{\strut\ignorespaces#1\strut}}
\newcommand{\CSLLeftMargin}[1]{\parbox[t]{\csllabelwidth}{\strut#1\strut}}
\newcommand{\CSLRightInline}[1]{\parbox[t]{\linewidth - \csllabelwidth}{\strut#1\strut}}
\newcommand{\CSLIndent}[1]{\hspace{\cslhangindent}#1}



\setlength{\emergencystretch}{3em} % prevent overfull lines

\providecommand{\tightlist}{%
  \setlength{\itemsep}{0pt}\setlength{\parskip}{0pt}}



 


\setlength\heavyrulewidth{0ex}
\setlength\lightrulewidth{0ex}
\usepackage[automark]{scrlayer-scrpage}
\clearpairofpagestyles
\cehead{
  Brian Weatherson
  }
\cohead{
  Précis of Knowledge: A Human Interest Story
  }
\ohead{\bfseries \pagemark}
\cfoot{}
\makeatletter
\newcommand*\NoIndentAfterEnv[1]{%
  \AfterEndEnvironment{#1}{\par\@afterindentfalse\@afterheading}}
\makeatother
\NoIndentAfterEnv{itemize}
\NoIndentAfterEnv{enumerate}
\NoIndentAfterEnv{description}
\NoIndentAfterEnv{quote}
\NoIndentAfterEnv{equation}
\NoIndentAfterEnv{longtable}
\NoIndentAfterEnv{abstract}
\renewenvironment{abstract}
 {\vspace{-1.25cm}
 \quotation\small\noindent\emph{Abstract}:}
 {\endquotation}
\newfontfamily\tfont{EB Garamond}
\addtokomafont{disposition}{\rmfamily}
\addtokomafont{descriptionlabel}{\rmfamily}
\addtokomafont{title}{\normalfont\itshape}
\let\footnoterule\relax

\makeatletter
\renewcommand{\@maketitle}{%
  \newpage
  \null
  \vskip 2em%
  \begin{center}%
  \let \footnote \thanks
    {\itshape\huge\@title \par}%
    \vskip 0.5em%  % Reduced from default
    {\large
      \lineskip 0.3em%  % Reduced from default 0.5em
      \begin{tabular}[t]{c}%
        \@author
      \end{tabular}\par}%
    \vskip 0.5em%  % Reduced from default
    {\large \@date}%
  \end{center}%
  \par
  }
\makeatother
\RequirePackage{lettrine}

\renewenvironment{abstract}
 {\quotation\small\noindent\emph{Abstract}:}
 {\endquotation\vspace{-0.02cm}}

\setmainfont{EB Garamond Math}[
  BoldFont = {EB Garamond SemiBold},
  ItalicFont = {EB Garamond Italic},
  RawFeature = {+smcp},
]

\newfontfamily\scfont{EB Garamond Regular}[RawFeature=+smcp]
\renewcommand{\textsc}[1]{{\scfont #1}}

\renewcommand{\LettrineTextFont}{\scfont}
\usepackage{amsfonts}
\usepackage{lettrine}
\KOMAoption{captions}{tableheading}
\makeatletter
\@ifpackageloaded{caption}{}{\usepackage{caption}}
\AtBeginDocument{%
\ifdefined\contentsname
  \renewcommand*\contentsname{Table of contents}
\else
  \newcommand\contentsname{Table of contents}
\fi
\ifdefined\listfigurename
  \renewcommand*\listfigurename{List of Figures}
\else
  \newcommand\listfigurename{List of Figures}
\fi
\ifdefined\listtablename
  \renewcommand*\listtablename{List of Tables}
\else
  \newcommand\listtablename{List of Tables}
\fi
\ifdefined\figurename
  \renewcommand*\figurename{Figure}
\else
  \newcommand\figurename{Figure}
\fi
\ifdefined\tablename
  \renewcommand*\tablename{Table}
\else
  \newcommand\tablename{Table}
\fi
}
\@ifpackageloaded{float}{}{\usepackage{float}}
\floatstyle{ruled}
\@ifundefined{c@chapter}{\newfloat{codelisting}{h}{lop}}{\newfloat{codelisting}{h}{lop}[chapter]}
\floatname{codelisting}{Listing}
\newcommand*\listoflistings{\listof{codelisting}{List of Listings}}
\makeatother
\makeatletter
\makeatother
\makeatletter
\@ifpackageloaded{caption}{}{\usepackage{caption}}
\@ifpackageloaded{subcaption}{}{\usepackage{subcaption}}
\makeatother
\usepackage{bookmark}
\IfFileExists{xurl.sty}{\usepackage{xurl}}{} % add URL line breaks if available
\urlstyle{same}
% fallback for those not using the hyperref driver hyperxmp:
\makeatletter
\@ifundefined{xmpquote}{\newcommand{\xmpquote}[1]{#1}}{}
\makeatother
\hypersetup{
  pdftitle={Précis of Knowledge: A Human Interest Story},
  pdfauthor={Brian Weatherson},
  pdfkeywords={\xmpquote{pragmatic
encroachment}, \xmpquote{interest-relativity}, \xmpquote{knowledge}, \xmpquote{rational
action}, \xmpquote{inquiry}, \xmpquote{epistemology}},
  hidelinks,
  pdfcreator={LaTeX via pandoc}}


\title{Précis of \emph{Knowledge: A Human Interest Story}}
\author{Brian Weatherson}
\date{2026}
\begin{document}
\maketitle
\begin{abstract}
In \emph{Knowledge: A Human Interest Story}, I argue for an
interest-relative theory of knowledge. The argument rests on principles
connecting knowledge and action, combined with some judgments about
which actions are rational in a particular low-stakes game. In this
short piece, I set out that argument, and highlight the distinctive
features of version of interest-relativity I defend: that odds rather
than stakes matter to knowledge; that merely hypothetical and
theoretical questions can change what one knows; that questions about
what to do must be separated from questions about what maximises
expected utility; and that knowledge is a constraint on the starting
points of inquiry, not its endings.
\end{abstract}


\setstretch{1.1}
My original plan for \emph{Knowledge: A Human Interest Story} was to
concatenate my various writings on interest-relativity over the years.
That plan failed, because it turned out I'd changed my mind on too many
things. The result is, I hope, better than the original plan; the
presentation of a coherent theory built out of the not always consistent
ideas in those earlier papers.

Like many people, I started thinking about interest-relativity because
of work by Jeremy Fantl and Matthew McGrath
(\citeproc{ref-FantlMcGrath2002}{2002}). In that paper, and in the
literature that followed, there were two big arguments for
interest-relativity. One was an argument from cases. Those arguments
made interest-relativity appear to be both similar to, and also a
competitor to, contextualism. The cases that motivated it were similar,
and a lot of the development of the argument was designed to show that
interest-relative theories provided a better explanation of the
intuitions about cases than contextualism.

I've long been sceptical of that line of argument, and instead more
interested in a different line that, I think, was always present. That
line starts with the clash between three principles and an intuition
about rational action. The three principles are:

\begin{enumerate}
\def\labelenumi{\arabic{enumi}.}
\tightlist
\item
  We know a lot.
\item
  What we know can be taken for granted in decision making.
\item
  What we know is independent of which decision we are facing.
\end{enumerate}

Then the intuition is that for some things we know, and some bets that
win if the known thing is true, it is irrational to take the bet. The
argument for interest-relativity is that something has to give, and the
least costly thing to give up is the decision-invariance of knowledge.
What we know depends on what decision we're making.

This argument doesn't rely on intuitions about what is or is not
knowledge. It does rely on intuitions about rational action, but I think
these are fairly widely shared. And it does rely on various principles,
which arguably rest in turn on theoretical intuitions. But it doesn't
have a role for intuitions about cases of knowledge.

The book is structured around a version of this argument, focussing on
the kind of bets involved in a game I call the Red-Blue game. Here's how
the game works. Two sentences are displayed to the player, one in red,
the other in blue. The player has to pick a colour and a truth-value.
They win iff the sentence with the colour they pick has the truth-value
they pick. So if they say ``Red-True'', they win iff the red sentence is
true. `Winning' here means getting a non-trivial but not huge reward; in
the book I make it \$50. The players know these rules of the game and
don't know anything else relevant about the game. Now consider an
instance of the game with these sentences.

\begin{description}
\tightlist
\item[Red]
Two plus two equals four.
\item[Blue]
The Battle of Agincourt was in 1415.
\end{description}

For people with normal levels of historical expertise, the only rational
play here is Red-True. Or, at least, that's the intuition I want to rest
on, and it's one that has been \emph{fairly} widely shared by people I
speak to. Now here's a sceptical argument. If the player knew that the
Battle of Agincourt was in 1415, they would know Blue-True is a winning
move, while they'd only know that Red-True wins if two plus two equals
four. So Blue-True would, in a sense, weakly dominate Red-True. So it's
at least \emph{permissible} to play Blue-True. But it's not permissible
to play Blue-True, so the player does not know that the Battle of
Agincourt was in 1415. By varying the blue sentence, we can similarly
argue that the player knows very little, contradicting the claim `we
know a lot'.

The solution I prefer is that once this game comes up, what the player
knows changes. They may have known when the Battle of Agincourt was, and
they may know it tomorrow, but they don't know it now. That's the
canonical instance of interest-relativity in the book. This is still, I
think, basically the same argument as in earlier arguments for
interest-relativity, but focussing on this case has several distinctive
features.

I've already mentioned one of these features. The book does not rest on
an intuition that the player loses knowledge. To be honest, even after
all these years, I still find that counter-intuitive. It rests on an
argument that the player loses knowledge, which in turn rests on
intuitions about rational play, and principles connecting knowledge and
rational play.

Another obvious consequence is that the view has nothing to do with high
\emph{stakes}. I didn't put more at stake than \$50. I don't think the
intuitions change much if you lower the payouts a fair bit from that. I
say, here agreeing with Schroeder (\citeproc{ref-Schroeder2012}{2012}),
that what matters are the odds the player faces, not the stakes.

What if the player did not in fact face this choice, but is idly
planning about what to do should it come up? This might strike you as a
little over the top, but I've spent quite a bit of time wondering about
this in the course of the book. I say the argument that they don't know
when the Battle was still goes through. Just being interested in a
question can change what one knows. That's because knowing something
requires being properly willing to treat it as well supported enough to
use it, should it be relevant, in any inquiry you are engaged
in.\footnote{The `well supported enough to use' locution does some work
  in the presentation in KAHIS; it's meant to be compatible with not
  using a piece of knowledge for reasons other than it's lack of
  evidential support. I now think it would be better to make the
  distinction I make in the Replies (later in this symposium) between
  finding out inquiries, and checking inquiries, and say that anything
  you know can, if relevant, be used in any finding out inquiry. But
  that's not what I said in KAHIS.}

In earlier work, I understood questions about what one should be willing
to use as being settled by the \emph{outputs} of the inquiry. The above
argument inferred that one should not start inquiry with knowledge of
the Battle's date from the fact that this would lead to the wrong
conclusion: Play Blue-True! I used to think that was what made it wrong
to use the Battle's date. This extra step is, I now think, a mistake.
Knowledge is not a constraint on where one ends inquiry, but on where
one starts it. The relevant constraint is that for someone inquiring
into \emph{Q?}, and \emph{p} is relevant in the right way to that
question, they know \emph{p} only if the questions \emph{Q?} and
\emph{If p, Q?} are in a deep sense the same question, and they are
answered in the same way. If \emph{p} is not something relevant to
anything one is currently inquiring into, the constraint is that there
is some possible inquiry where one would use \emph{p} as the first step.
This move avoids some very implausible consequences of what I'd
previously said about knowledge of practically (and theoretically)
irrelevant propositions.

I've spoken liberally about questions here, but the last few points that
are distinctive to the book turn on being a bit more careful about what
a question is. In earlier work I'd equated, sometimes explicitly,
sometimes implicitly, the following two questions:

\begin{itemize}
\tightlist
\item
  Which live option maximises expected utility?
\item
  What should I do?
\end{itemize}

But these are not the same question, and some of the challenges for an
interest-relative view turn on getting clear on the difference between
them. Here are four important differences.

First, if two options have the same expected utility, I still might have
to decide which to do. I want a can of Coke, there are two on the shelf.
It would do just as well to take either, but I have to actually take
one. Knowing the answer to the utility question, they have the same
utility, doesn't settle what my arm should do.

Second, if two options have \emph{really similar} expected utility, I
shouldn't always take the one with higher expected utility. It might be
more trouble than it's worth to calculate which has higher utility, and
I might be (faultlessly) unreliable at choosing the higher expected
utility without calculating. Maybe it would be less work by some
miniscule amount to pick up the left can rather than the right can. It
could, I say, still be practically rational to act like the cans are
equally good, and choose arbitrarily. In chapter 6, I show that some
objections to interest-relative accounts, which work as long as we
equate these questions, no longer work if we separate them.

Third, it's not obvious how to use the expected utility framework to
model mathematical uncertainty. But mathematical uncertainty might be
something we want the interest-relative story to capture. If I'm playing
the red-blue game, and the blue sentence is \emph{67 plus 58 equals
125}, my intuition is that only Red-True is rational to choose. I'm just
more likely to make an error about this calculation than about two plus
two. But working this into an expected utility model raises extra
complications.

Fourth, an expected utility comparison needs a probability function, and
the usual picture is that the probability in question is
\emph{evidential} probability. This raises problems if there are
propositions that, if known, are part of one's evidence, and whether
they are known is interest-relative. For instance, imagine that there is
a Ford outside my window, in fact a Mustang with a distinctive shape,
and I'm looking straight at it. I'm asked to play the red-blue game with
the blue sentence being \emph{There is a Ford outside my window}. I
think in that context I don't know there is a Ford outside, so it's not
part of my evidence there is a Ford outside. This makes the very notion
of expected utility complicated. In chapter 8, I suggest a way of
dealing with that complication.

The last new point in the book which I'll mention here concerns
double-checking. The argument from the red-blue game to
interest-relativity requires some principles linking knowledge and
action. It has been argued, quite plausibly, that these principles are
inconsistent with the rationality of double checking something one
already knows. After all, if one knows \emph{p}, then checking \emph{p}
seems worse than just reasoning from one's knowledge that \emph{p} must
be true. I argue, in chapter 5, this isn't right. There are many goals
to inquiry. One possible goal is forming a more \emph{sensitive} belief
about some subject. Even if one knows \emph{p}, then an inquiry that
aims to maximise sensitivity should not simply reason from \emph{p} to
the truth of \emph{p}. This is, I hope, an independently interesting
claim about inquiry, but the role it plays here is to defend principles
connecting knowledge and action from counterexamples involving
double-checking.

\section*{References}\label{references}
\addcontentsline{toc}{section}{References}

\protect\phantomsection\label{refs}
\begin{CSLReferences}{1}{0}
\bibitem[\citeproctext]{ref-FantlMcGrath2002}
Fantl, Jeremy, and Matthew McGrath. 2002. {``Evidence, Pragmatics, and
Justification.''} \emph{Philosophical Review} 111 (1): 67--94. doi:
\href{https://doi.org/10.2307/3182570}{10.2307/3182570}.

\bibitem[\citeproctext]{ref-Schroeder2012}
Schroeder, Mark. 2012. {``Stakes, Withholding and Pragmatic Encroachment
on Knowledge.''} \emph{Philosophical Studies} 160 (2): 265--85. doi:
\href{https://doi.org/10.1007/s11098-011-9718-1}{10.1007/s11098-011-9718-1}.

\end{CSLReferences}



\noindent Draft.


\end{document}
